\section{Огляд літератури}
\label{sec:related}

\emph{BFT-протоколи.}
Класичний PBFT~\cite{castro1999} допускає $f<\n/3$ при $O(\n^{2})$
повідомленнях~\cite{lamport1982,dwork1988}. Пізніші конструкції ---
HotStuff~\cite{yin2019}, Tendermint~\cite{buchman2016},
Algorand~\cite{gilad2017}, HoneyBadgerBFT~\cite{miller2016} --- обмінюють
затримку та рандомізацію; Vukoli\'{c}~\cite{vukolic2015} стверджує, що жоден
протокол не домінує. Ці лінії робіт фіксують множину валідаторів; вони не
визначають її розмір за сумісними вимогами до пропускної здатності та детекції.

\emph{Бенчмарки.}
Fabric~\cite{androulaki2018,thakkar2018}, BLOCKBENCH~\cite{dinh2017} та
BFT-SMaRt~\cite{sousa2018} звітують абсолютну пропускну здатність у вибраних
конфігураціях. Вони не відокремлюють конкурентність від кількості валідаторів і
не формулюють умов вимірювання, за яких очікується перенесення підігнаних
показників.

\emph{Характеризація масштабованості та інструменти.}
Amdahl~\cite{amdahl1967}, Gustafson~\cite{gustafson1988} і USL
Gunther'а~\cite{gunther2007} згортають конкурентність і розмір системи в одну
змінну навантаження --- змішування, яке \RNM{} покликаний усунути. Межі
концентрації~\cite{hoeffding1963,boucheron2013}, design effect
Kish'а~\cite{kish1965}, інтервали методу дельти~\cite{seber2003}, залишковий
бутстреп~\cite{efron1993}, зміщення через пропущену змінну в неправильно
специфікованих регресіях~\cite{wooldridge2010} та програмування з імовірнісними
обмеженнями~\cite{charnes1959,nemirovski2007} --- стандартні інструменти.

\emph{Попередні роботи.}
Аналізи дистанційного голосування~\cite{springall2014,specter2020,park2021}
виносять загрози кінцевих точок і примусу поза область. Наше попереднє
калібрування~\cite{peniak2026a} і дослідження перемикання~\cite{peniak2026b}
залишили відкритими ідентифікований показник консенсусу та визначення розміру з
виміряними частотами детектора.
